\documentclass[11pt]{article}
\usepackage[T1]{fontenc}
\usepackage{lmodern}
\usepackage[margin=1in]{geometry}
\usepackage{amsmath,amssymb,amsthm,mathtools,microtype,booktabs}
\usepackage{xcolor,xurl}
\usepackage[colorlinks=true,linkcolor=blue!45!black,citecolor=blue!45!black,urlcolor=blue!45!black]{hyperref}
\hypersetup{pdftitle={Transcendence targets for the friend-tree exponent},pdfauthor={Christopher D. Long}}
\setlength{\emergencystretch}{2em}
\allowdisplaybreaks
\newtheorem{theorem}{Theorem}[section]
\newtheorem{proposition}[theorem]{Proposition}
\newtheorem{corollary}[theorem]{Corollary}
\newtheorem{lemma}[theorem]{Lemma}
\theoremstyle{remark}
\newtheorem{remark}[theorem]{Remark}
\newcommand{\Q}{\mathbb Q}
\newcommand{\C}{\mathbb C}
\newcommand{\E}{\mathbb E}
\newcommand{\Z}{\mathbb Z}
\newcommand{\OO}{\mathcal O}
\DeclareMathOperator{\den}{den}
\DeclareMathOperator{\lcm}{lcm}
\DeclareMathOperator{\Norm}{N}
\title{Transcendence targets for the friend-tree exponent:\\
exact first-promotion reduction and arithmetic obstructions}
\author{Christopher D. Long}
\date{November 23, 2026}
\begin{document}
\maketitle
\begin{abstract}
We investigate the arithmetic of the original friend-tree growth exponent
through an exact component of its inert-parent kernel. Eliminating all
null rings before the first promotion produces an explicit ${}_3F_2$
transform. We reduce the whole waiting phase to a positive residual with
a coefficient that is nonzero at every positive discount, classify its
rational telescopes, and give accelerated rational enclosures. For rational
discount $s=p/q>0$, the primitive coefficient denominators and their least
common denominator have logarithm $n\log n+O_s(n)$ when $p^2+4q^2$ is
nonsquare; the square case has exponential common denominators. A
prime-ideal calculation gives an exact coefficient-height asymptotic at
every positive algebraic discount. These statements identify the arithmetic
class of this genuine first-promotion kernel. We also record the scope of
Sprind\v zuk's applicable $E^*$-function theorem and an exact equivalence
for the local-model transcendence problem. Neither irrationality nor
transcendence of the original exponent is proved. The remaining step is
the construction of a nonzero signed critical witness with an arithmetic
error estimate strong enough to beat algebraic separation.
\end{abstract}

\section{Scope and the critical arithmetic question}

The original random friend tree selects a uniformly random existing
vertex, selects one of its neighbours uniformly, and attaches a new
vertex to that neighbour\cite{KR}. Let $\mu$ denote its core-growth
exponent. The present project\cite{Long}, using the upstream existence
theorem and certificates\cite{Uzay}, has the exact characterization
\[
 r_\infty(\mu)=1,\qquad
 I:=\left[\frac{113}{200},\frac{2832242711}{5000000000}\right]
 \ni\mu.
\]
Here $r_\infty$ is the inert-newborn generation pressure. It retains the
evolving environment; no finite independent-family characteristic
equation is asserted for this process. The bounds, existence theorem,
and exact pressure characterization are contextual inputs, not results
reproved in this note.

The new computations below concern the actual \emph{first proper
promotion} from the inert newborn. This state has one real core, one
real leaf, and one inert parent stub. The stub counts toward degree but
has no source clock. It differs from a free star, whose first core ring
always promotes a leaf. A first-promotion transform is an exact kernel
component; its evaluation is not a new equation for $\mu$.

\section{The exact first-promotion transform}

Let $\tau$ be the first promotion time from this inert newborn, and let
$J$ be the number of leaves immediately before that promotion. Define
\[
 \Phi(s)=\E e^{-s\tau},\qquad s>0.
\]
At leaf count $j\ge1$, leaf addition has rate $j$, promotion has rate
$j/(j+1)$, and the root's choice of its stub is a null event of rate
$1/(j+1)$. The promotion hazard is at least $1/2$, so $\tau$ is finite
almost surely. Exact geometric elimination of null events gives proper
holding rate $j+j/(j+1)$.

\begin{theorem}[Whole first-promotion kernel]\label{thm:kernel}
Put
\[
 a=\frac{s+2+\sqrt{s^2+4}}2,\qquad
 b=\frac{s+2-\sqrt{s^2+4}}2.
\]
Then $a>2$, $0<b<1$, $a+b=s+2$, and $ab=s$. With $t_0(s)=1$, set
\begin{equation}\label{eq:tn}
 t_n(s)=\prod_{j=1}^n\frac{j^2}{j^2+(s+2)j+s}
       =\frac{(n!)^2}{(a+1)_n(b+1)_n}.
\end{equation}
For $n\ge1$,
\begin{equation}\label{eq:joint}
 t_n(s)=\E[e^{-s\tau};J=n],\qquad
 \Phi(s)=\sum_{n\ge1}t_n(s)
 ={}_3F_2\!\left(\begin{matrix}1,1,1\\a+1,b+1\end{matrix};1\right)-1.
\end{equation}
Moreover
\begin{equation}\label{eq:asymp}
 t_n(s)\sim\Gamma(a+1)\Gamma(b+1)n^{-s-2}.
\end{equation}
\end{theorem}
\begin{proof}
At leaf count $j$, the discounted continuation and promotion factors are
respectively
\[
 \frac{j(j+1)}{j^2+(s+2)j+s},\qquad
 \frac{j}{j^2+(s+2)j+s}.
\]
Multiply the continuation factors for $1\le j<n$ and the final
promotion factor at $n$. Their numerator is $(n!)^2$, giving
\eqref{eq:joint}. Summing disjoint promotion outcomes gives $\Phi$.
The hypergeometric notation follows directly from its defining
series\cite{DLMF}; the excess of denominator parameters over numerator
parameters is $s+1>0$. Formula \eqref{eq:asymp} follows from the
gamma-ratio expression for the rising factorials.
\end{proof}

At $s=0$, the limiting formula is $t_n(0)=2/[(n+1)(n+2)]$, with
$\Phi(0)=1$. For positive $s$ there is also the positive double-integral
representation
\begin{equation}\label{eq:double}
 \Phi(s)=s\int_0^1\int_0^1
 \frac{xy}{1-xy}(1-x)^{a-1}(1-y)^{b-1}\,dx\,dy.
\end{equation}
Indeed each factor $n!/(a+1)_n$ is $a\,B(n+1,a)$; Tonelli's theorem
then sums the geometric series. Thus $(t_n)_{n\ge0}$ is exactly the
Hausdorff moment sequence of the product of independent
$\operatorname{Beta}(1,a)$ and $\operatorname{Beta}(1,b)$ variables.
No moment representation for the full friend-tree mean sequence is
inferred from this statement.

\section{Reduction, nonvanishing, and boundary terms}

Put
\[
 D_s(z)=(z+a+1)(z+b+1)=z^2+(s+4)z+2s+3,
\]
and define the rational difference operator
\begin{equation}\label{eq:delta}
 (\delta_sR)(z)=R(z)-\frac{(z+1)^2}{D_s(z)}R(z+1).
\end{equation}
It is the summation operator associated with $t_{n+1}/t_n$.

\begin{theorem}[Exact positive residual]\label{thm:reduction}
For every $s>0$,
\begin{equation}\label{eq:reduction}
 \boxed{\displaystyle
 \Phi(s)=\frac{2}{(s+1)(s+2)}+
 \frac{s^2}{(s+1)(s+2)}\sum_{n\ge1}\frac{t_n(s)}{n^2}.}
\end{equation}
The residual is strictly positive and its displayed coefficient does
not vanish at any positive specialization.
\end{theorem}
\begin{proof}
Let
\[
 R_s(z)=\frac{z}{s+1}+\frac{s^2+3s+4}{(s+1)(s+2)},
 \qquad c_s=\frac{s^2}{(s+1)(s+2)}.
\]
Direct polynomial multiplication gives
$\delta_sR_s=1-c_s/D_s$. Since
$t_n/D_s(n)=t_{n+1}/(n+1)^2$, summing the identity gives the exact finite
formula
\begin{equation}\label{eq:finite-boundary}
 \sum_{n=1}^N t_n
 =\frac{2}{(s+1)(s+2)}+c_s\sum_{n=1}^N\frac{t_n}{n^2}
  -t_N\bigl(R_s(N)-1\bigr).
\end{equation}
The final term tends to zero by \eqref{eq:asymp}. This proves the
identity and positivity.
\end{proof}

\begin{theorem}[Classification of rational telescopes]\label{thm:nosum}
For $s>0$,
\[
 1\in\delta_s\C(z)\quad\Longleftrightarrow\quad s\in\Z_{>0}.
\]
Consequently no rational telescope exists for any $s\in I$, whether
that parameter is rational, algebraic irrational, or transcendental.
\end{theorem}
\begin{proof}
A rational solution would satisfy
\begin{equation}\label{eq:poly-telescope}
 D_s(z)R(z)-(z+1)^2R(z+1)=D_s(z).
\end{equation}
If $R$ had finite poles, consider one of their orbits under integer
translation. At one step to the left of its leftmost pole, the second
term alone can have a pole. Its multiplier forces that leftmost pole
to be $0$. The rightmost pole in the same orbit is then a nonnegative
integer. At that point only the first term can have a pole, whereas
$D_s$ has both roots strictly negative. This is impossible. Hence $R$
is polynomial.

If its degree is $m$ and leading coefficient is $c$, the coefficient of
$z^{m+1}$ on the left of \eqref{eq:poly-telescope} is $c(s+2-m)$.
For $m\ge2$ it must vanish, so $m=s+2$. For $m\le1$, write $R=Az+B$.
Matching the quadratic and linear terms gives
\[
 A=\frac1{s+1},\qquad
 B=\frac{s^2+3s+4}{(s+1)(s+2)}.
\]
The remaining constant equation is $s^2=0$, excluded here. Thus any
positive parameter permitting a telescope is an integer.

Conversely, let $s=k$ be a positive integer. The operator
$R\mapsto D_kR-(z+1)^2R(z+1)$ maps polynomials of degree at most
$k+2$ to the same space, because the highest coefficient cancels.
It has trivial kernel: a homogeneous solution satisfies $R(-1)=0$,
then $R(-2)=R(-3)=\cdots=0$. In this induction $D_k$ is nonzero at
every integer, since $\sqrt{k^2+4}$ is irrational for $k\ge1$.
Hence the operator is invertible on this finite-dimensional space,
and \eqref{eq:poly-telescope} has a unique polynomial solution.
\end{proof}

\paragraph{Formal nonvanishing has a precise scope.}
The theorem excludes a rational antidifference. It does not show that
the evaluated sum is irrational, and it does not show that this class
survives a signed reduction of the global critical equation.

\begin{proposition}[A nonzero boundary at an exceptional parameter]
At $s=1$,
\begin{equation}\label{eq:s1}
 \Phi(1)=\Gamma\!\left(\frac{5+\sqrt5}{2}\right)
              \Gamma\!\left(\frac{5-\sqrt5}{2}\right)-3
        =-\frac{\pi}{\cos(\pi\sqrt5/2)}-3.
\end{equation}
\end{proposition}
\begin{proof}
The polynomial $R(z)=-z^3-5z^2-7z-2$ satisfies $\delta_1R=1$.
Since $t_n\sim\Gamma(a+1)\Gamma(b+1)n^{-3}$, its telescoping boundary
is $-\Gamma(a+1)\Gamma(b+1)$ rather than zero. Thus
$\sum_{n\ge0}t_n=\Gamma(a+1)\Gamma(b+1)-2$.
Now $a+b=3$, $ab=1$, and $(a-1)(a-2)=1$; the gamma reflection
formula gives the last equality in \eqref{eq:s1}.
\end{proof}
This is an exact evaluation, not an arithmetic classification. The
transcendence of an individual factor does not classify its quotient
with $\pi$.

\section{Accelerated rational enclosures}

The positive reduction also yields a practical finite calculator for
this kernel component. Define
\[
 u_n=t_n/n^2,\quad A=s+4,\quad B=2s+3,\quad \kappa=s+3.
\]
Then $u_{n+1}/u_n=n^2/(n^2+An+B)$.

\begin{theorem}[Two rational tail bounds]\label{thm:tail}
For an integer $M\ge1$, put
\[
 c_-=\kappa-B/A,\qquad
 c_+(M)=\kappa-\frac{BM}{AM+B}.
\]
Then
\begin{equation}\label{eq:tail}
 L_M:=\frac{u_M(M+c_-)}{\kappa}
 <\sum_{n\ge M}u_n<
 U_M:=\frac{u_M(M+c_+(M))}{\kappa},
\end{equation}
and
\begin{equation}\label{eq:tail-width}
 U_M-L_M=\frac{u_MB^2}{\kappa A(AM+B)}.
\end{equation}
All endpoints are rational at rational $s>0$.
\end{theorem}
\begin{proof}
For $R_c(n)=(n+c)/\kappa$ the residual is
\[
 R_c(n)-\frac{n^2R_c(n+1)}{n^2+An+B}-1
 =\frac{[A(c-\kappa)+B]n+B(c-\kappa)}
        {\kappa(n^2+An+B)}.
\]
For $c=c_-$ its numerator is $-B^2/A<0$. For $c=c_+(M)$ its
numerator is $B^2(n-M)/(AM+B)\ge0$ for $n\ge M$, strictly positive
when $n>M$. Multiplying by $u_n$, summing, and using $u_nR_c(n)\to0$
proves \eqref{eq:tail}. Subtracting gives \eqref{eq:tail-width}.
\end{proof}

Insert these endpoints into \eqref{eq:reduction} after summing
$u_1,\ldots,u_{M-1}$. The resulting interval for $\Phi(s)$ has exact
width
\begin{equation}\label{eq:phi-width}
 \frac{s^2}{(s+1)(s+2)}\frac{t_M}{M^2}
 \frac{(2s+3)^2}{(s+3)(s+4)((s+4)M+2s+3)}
 =O_s(M^{-s-5}).
\end{equation}
The unaccelerated omitted sum is instead of order $M^{-s-1}$.
For example, the included exact checker with $M=51$ gives
\[
 0.520122548676957758785996459558
 <\Phi(13/23)<
 0.520122548689960716256732865482.
\]
These decimal endpoints are rounded outwards from exact fractions.
They are not new estimates of $\mu$.

\section{Primitive denominator growth at rational discounts}

Let $s=p/q>0$ be reduced, with $p,q$ positive integers. Set
\begin{equation}\label{eq:quadratic}
 Q_s(j)=qj^2+(p+2q)j+p,\qquad \Delta=p^2+4q^2.
\end{equation}
Then
\[
 t_n=\frac{q^n(n!)^2}{\prod_{j=1}^n Q_s(j)}.
\]
Write $t_n=A_n/B_n$ in lowest positive terms and
$\mathcal D_n=\lcm(B_0,\ldots,B_n)$.

\begin{theorem}[Exact denominator dichotomy]\label{thm:den}
If $\Delta$ is nonsquare, then
\begin{equation}\label{eq:den-asymp}
 \log A_n=n\log n+O_s(n),\qquad
 \log B_n=n\log n+O_s(n),\qquad
 \log\mathcal D_n=n\log n+O_s(n).
\end{equation}
If $\Delta$ is square, then
\begin{equation}\label{eq:square-den}
 \log\mathcal D_n=O_s(n).
\end{equation}
\end{theorem}
\begin{proof}
Call a prime $\ell$ good if $\ell\nmid2q\Delta$. At an inert good
prime, the quadratic $Q_s$ has no root modulo $\ell$, so
\[
 v_\ell(A_n)=2v_\ell(n!).
\]
At a split good prime, $Q_s$ has exactly two simple roots modulo every
$\ell^h$. Every complete block of length $\ell^h$ contains both
roots. Therefore
\[
 v_\ell\left(\prod_{j=1}^n Q_s(j)\right)
 \ge2\sum_{h\ge1}\left\lfloor n/\ell^h\right\rfloor
 =2v_\ell(n!),
\]
so $v_\ell(A_n)=0$. The contribution of all bad primes to
$\log A_n$ is $O_s(n)$, since $A_n$ divides $q^n(n!)^2$.

For the nontrivial quadratic character attached to
$\Q(\sqrt\Delta)$, the classical weighted Mertens estimate gives
\[
 \sum_{\substack{\ell\le n\\\ell\text{ good inert}}}
      \frac{\log\ell}{\ell-1}
 =\tfrac12\log n+O_s(1).
\]
This uses only the usual prime distribution in a fixed quadratic
field, with no RH assumption; see the unconditional character-sum
form recorded in\cite[Section 1.2]{Zhao}. Legendre's formula and the
bound $\pi(n)=O(n/\log n)$ now give
\[
 \log A_n=2\sum_{\ell\text{ good inert}}v_\ell(n!)\log\ell+O_s(n)
 =n\log n+O_s(n).
\]
Since $\log t_n=-(s+2)\log n+O_s(1)$, the same leading asymptotic
holds for $B_n$.

To control the least common denominator, first consider a good split
prime $\ell\le n$. For each $h$, the excess root count over
$2\lfloor k/\ell^h\rfloor$ belongs to $\{0,1,2\}$. Since
$Q_s(j)\le C_s n^2$ for $j\le n$, it follows that
\[
 v_\ell(B_k)\le2\log_\ell(C_s n^2)\qquad(k\le n).
\]
Summing this bound over primes at most $n$ costs $O_s(n)$ in the
logarithm. Good inert primes contribute no denominators. For a fixed
bad prime, let $\widetilde A_k=q^k(k!)^2$ and
$\widetilde B_k=\prod_{j\le k}Q_s(j)$. The latter products are
nested, so the excess of $\max_{k\le n}v_\ell(B_k)$ over
$v_\ell(B_n)$ is at most $v_\ell(\widetilde A_n)=O_s(n)$.
For primes $\ell>n$, once $n$ exceeds the fixed bad primes, the
numerator has no $\ell$ factor and denominator valuations increase
with $k$. Hence $B_n$ already attains their maxima. Altogether
$0\le\log\mathcal D_n-\log B_n=O_s(n)$.

If $\Delta$ is square, $a,b$ are positive rational numbers. For
$a=u/v$ in lowest terms, the common denominator of
$k!/(a+1)_k$, $0\le k\le n$, divides
$\lcm(1,2,\ldots,u+vn)$. Indeed, comparison with Legendre's formula
leaves at most one extra valuation for each prime power. This lcm has
logarithm $O_a(n)$. Apply the same argument to $b$ and multiply the
two clearing factors.
\end{proof}

The exceptional rational parameters are dense: writing
\[
 s=t-t^{-1},\qquad t\in\Q,\quad t>1,
\]
gives $\sqrt{s^2+4}=t+t^{-1}$. They cannot be excluded using the
current enclosure. For example
\[
 \frac{2603}{4602}\in I,\qquad
 \sqrt{(2603/4602)^2+4}=\frac{9565}{4602}.
\]
The nonsummability theorem nevertheless applies at these parameters,
since none in $I$ is an integer.

\section{An exact height law at every algebraic discount}

The primitive arithmetic admits a stronger all-degree formulation.
We use the absolute logarithmic Weil height $h$.

\begin{theorem}[Exact height of the first-promotion coefficients]
\label{thm:alg-den}
Let $s>0$ be algebraic, and define
\[
 \vartheta=\frac{s+\sqrt{s^2+4}}2>1,\qquad
 K=\Q(s),\quad L=\Q(\vartheta),\quad
 d=[K:\Q],\quad e=[L:K]\in\{1,2\},\quad D=ed.
\]
Then $s=\vartheta-\vartheta^{-1}$, and
\begin{equation}\label{eq:alg-den}
 \boxed{\displaystyle
 h(t_n(s))=\left(2-\frac2D\right)n\log n+O_s(n).}
\end{equation}
More precisely, writing $(t_n)=\mathfrak a_n\mathfrak b_n^{-1}$
with coprime integral ideals of $K$, both
$\log\Norm\mathfrak a_n$ and $\log\Norm\mathfrak b_n$ equal
\[
 \left(2d-\frac2e\right)n\log n+O_s(n),
\]
and their difference is $O_s(\log n)$.
The leading height coefficient vanishes exactly when $\vartheta$
is rational, equivalently when $s$ is rational and $s^2+4$ is a
rational square.
\end{theorem}
\begin{proof}
Write $q_s(j)=j(j+2)+(j+1)s$. Exclude a fixed finite set of rational
primes containing the coefficient denominators, the relevant order
indices and ramified primes, and primes dividing the relative
quadratic discriminant. At a remaining prime $p$, reduction of $s$
generates the residue field at each prime $\mathfrak p\mid p$ of $K$.
If $q_s(j)$ vanishes there, then $j\not\equiv-1\pmod p$ and
\[
 s\equiv-\frac{j(j+2)}{j+1}\in\mathbb F_p.
\]
Thus $\mathfrak p$ has residue degree one. At a degree-one prime
the quadratic has either zero or two distinct roots. In the two-root
case, Hensel lifting and complete residue blocks show that
\[
 v_{\mathfrak p}\!\left(\prod_{j=1}^nq_s(j)\right)
 \ge2v_p(n!),
\]
so there is no numerator-ideal contribution. At every other prime
of $K$ above $p$, the numerator valuation is exactly $2v_p(n!)$.

Let $F_p$ be the number of degree-one primes of $L$ over $p$.
The number of degree-one primes of $K$ where the quadratic splits
is $F_p/e$. For $e=2$ each splitting prime gives two such primes
of $L$; for $e=1$ the roots already lie in $K=L$. Since
\[
 a=1+\vartheta,\qquad b=1-\vartheta^{-1},
\]
the quadratic splitting field over $K$ is precisely $L$. Taking
ideal norms gives the exact good-prime identity
\begin{equation}\label{eq:ideal-valuations}
 v_p(\Norm\mathfrak a_n)
 =2\left(d-\frac{F_p}{e}\right)v_p(n!).
\end{equation}
At a fixed excluded prime, the bound
$v_{\mathfrak p}(q_s(j))\ge\min(0,v_{\mathfrak p}(s))$
shows that the positive valuation of $t_n$ is $O_s(n)$.
Consequently all excluded primes contribute only $O_s(n)$ to the
logarithm of the numerator norm.

The unconditional number-field Mertens estimate\cite[Theorem 1,
\textnormal{(A1)--(A2)}]{GarciaLee} is
\[
 \sum_{\Norm\mathfrak P\le x}
 \frac{\log\Norm\mathfrak P}{\Norm\mathfrak P}
 =\log x+O_L(1).
\]
For $L=\Q$ use classical Mertens. Prime ideals of residue degree
at least two have bounded total contribution to this sum. Replacing
$1/p$ by $1/(p-1)$ also changes it by a bounded amount, so
\[
 \sum_{p\le x}F_p\frac{\log p}{p-1}=\log x+O_L(1).
\]
Since $F_p\le D$ and
$v_p(n!)=n/(p-1)+O(\log n/\log p)$, it follows that
\[
 \sum_{p\le n}F_pv_p(n!)\log p=n\log n+O_L(n).
\]
Summing \eqref{eq:ideal-valuations} therefore yields
\[
 \log\Norm\mathfrak a_n
 =2d\log(n!)-\frac2e n\log n+O_s(n)
 =\left(2d-\frac2e\right)n\log n+O_s(n).
\]

For each embedding $\sigma:K\hookrightarrow\C$, no factor
$q_{\sigma(s)}(j)$ vanishes: such vanishing would make $\sigma(s)$
rational, impossible if $s$ is irrational, and contradict positivity
if $s$ is rational. The elementary expansion
\[
 \log|q_{\sigma(s)}(j)|
 =2\log j+\frac{\Re(\sigma(s)+2)}j+O_s(j^{-2})
\]
gives $\log|\sigma(t_n)|=O_s(\log n)$. Hence
$\log|\Norm(t_n)|=O_s(\log n)$, proving the same main term for
the denominator ideal. Finally
\[
 h(t_n)=\frac1d\left(\log\Norm\mathfrak b_n+
       \sum_{\sigma:K\hookrightarrow\C}\log\max(1,|\sigma(t_n)|)\right)
\]
proves \eqref{eq:alg-den}, since the archimedean sum is $O_s(\log n)$.
\end{proof}

\begin{corollary}[The arithmetic function class is exactly determined]
\label{cor:classes}
For positive algebraic $s$, put
\[
 H_s(z)=\sum_{n\ge0}t_n(s)z^n,\qquad
 E_s(z)=\sum_{n\ge0}t_n(s)z^n/n!.
\]
Then $H_s$ is a $G$-function, and $E_s$ an $E$-function, if and only
if $s$ is rational and $s^2+4$ is a rational square.
\end{corollary}
\begin{proof}
With $\theta=z\,d/dz$, the coefficient recurrence gives
\begin{align*}
 [\theta^2+(s+2)\theta+s-z(\theta+1)^2]H_s&=s,\\
 [\theta^2+(s+2)\theta+s-z(\theta+1)]E_s&=s.
\end{align*}
Differentiation homogenizes these equations, so both functions are
$D$-finite. Their coefficient conjugates have polynomial size. The
square rational case therefore satisfies all the usual exponential
arithmetic conditions. In every other algebraic case,
Theorem~\ref{thm:den} or \ref{thm:alg-den} rules those conditions out,
even for an individual coefficient. In particular the obstruction also
violates Siegel's weaker subfactorial condition\cite{Beukers}.
\end{proof}

This classification concerns these specific generating functions, not
every possible representation of $\Phi(s)$. Moreover
$\Phi(s)=H_s(1)-1$; even in the exceptional $G$-function case, the
$E$-function $E_s$ has not been substituted for this value. A Borel
transform does not preserve the evaluation required by the tree.

\section{The local-model specialization problem and existing theory}

The companion local models have a different, simpler exact transform
\[
 F(x)=\int_0^1u^x e^{u-1}\,du.
\]
Their root equation is
\begin{equation}\label{eq:van}
 F(\alpha_\omega)=
 \frac{(1-\alpha_\omega)(\alpha_\omega+\omega)}
      {\alpha_\omega+(1-\alpha_\omega)\omega}.
\end{equation}
The project proves irrationality for every finite algebraic
$\omega\ge0$ via rational specialization and Siegel--Shidlovskii
theory\cite{Long,Beukers}. That argument does not already prove
transcendence.

\begin{proposition}[An exact equivalence for the entire local family]
Let $\alpha_0$ be the PAN root. The following assertions are equivalent:
\begin{enumerate}
\item $\alpha_\omega$ is transcendental for every finite algebraic
      $\omega\ge0$.
\item $F(x)$ is transcendental for every algebraic real
      $x\in[\alpha_0,1)$.
\end{enumerate}
\end{proposition}
\begin{proof}
Solving \eqref{eq:van} for the activation parameter gives
\[
 \Omega(x)=\frac{x(F(x)+x-1)}{(1-x)(1-F(x))},\qquad
 \Q(x,F(x))=\Q(x,\Omega(x)).
\]
Here $0<F(x)<1$ for $0<x<1$. The function $F(x)+x-1$ is strictly
increasing since its derivative is at least $1-(x+1)^{-2}>0$.
It vanishes at $\alpha_0$, so $\Omega(x)\ge0$ on the stated interval.
Uniqueness of the local root gives $\alpha_{\Omega(x)}=x$.
Thus an algebraic value $F(x)$ at an algebraic $x$ gives an admissible
algebraic $\Omega(x)$ with algebraic exponent. Conversely every
algebraic local exponent gives an algebraic $F$-value through
\eqref{eq:van}. The same rearrangement shows every $\alpha_\omega$
with $\omega\ge0$ lies in $[\alpha_0,1)$.
\end{proof}

There is relevant arithmetic theory beyond ordinary $E$-functions.
Sprind\v zuk\cite[Theorems 1 and 3]{Sprindzuk} considers precisely
\[
 f_\lambda(z)=\sum_{n\ge0}\frac{z^n}{(\lambda+1)_n},\qquad
 zf_\lambda'(z)=(z-\lambda)f_\lambda(z)+\lambda.
\]
For fixed algebraic irrational $\lambda$ and positive integers $D,t$,
only finitely many algebraic $z$ of degree at most $t$ have
$f_\lambda(z)$ algebraic of degree at most $D$. His theorem applies
to this function through the broader $E^*$ class. Since
\[
 F(x)=f_{x+1}(-1)/(x+1),
\]
it addresses the correct family. It does not exclude the particular
point $-1$, and its finite exceptional set changes with $\lambda$.
It therefore gives no finite exceptional list of candidate exponents.
Removing the possible exceptional value at $-1$ is a precise arithmetic
target for the local model. No implication for the original $\mu$ is
available without a separate exact transfer.

\section{The next critical witness and a dependency audit}

The new reduction proves that a complete first-promotion phase contains
a nonzero rational summation class throughout $I$. The denominator
theorems explain why positive truncation and fixed-order acceleration
are insufficient for the intended small-integer argument. They do not
rule out an improved signed construction or cancellations across the
global kernel.

One precise transcendence target is to construct integer polynomials
$P_N$ from the actual critical equation such that, infinitely often,
$P_N(\mu)\ne0$, and
\begin{equation}\label{eq:target}
 \frac{-\log|P_N(\mu)|}
 {\log(2H(P_N))+\deg P_N+\log(\deg P_N+2)}\longrightarrow\infty.
\end{equation}
Both nonvanishing and the quantitative estimate are unproved here.
This is a sufficient criterion, not a claim that every transcendence
proof must have this form. To see why it suffices, suppose $\mu$ were
algebraic of degree $d$, with primitive minimal polynomial of leading
coefficient $c$ and conjugates $\mu_1=\mu,\ldots,\mu_d$. The nonzero
integer resultant gives, for $m=\deg P$,
\[
 1\le |c|^m|P(\mu)|\prod_{j=2}^d|P(\mu_j)|
 \le |P(\mu)|[(m+1)H(P)]^{d-1}
       \left(|c|\prod_{j=2}^d\max(1,|\mu_j|)\right)^m.
\]
This contradicts \eqref{eq:target}. For a fixed proposed algebraic
number, a sufficiently strong fixed rate would already suffice.

The substantive missing construction is a \emph{signed critical
boundary functional} that both detects $r_\infty(s)-1$ and preserves a
nonzero arithmetic class after the full environment is incorporated.
The identity for $\Phi$ does not construct that functional. In
particular $\Phi(\mu)=1$ is false: $\Phi(s)<1$ for every $s>0$.

\begin{center}\small
\begin{tabular}{p{0.27\textwidth}p{0.62\textwidth}}
\toprule
Statement & Dependency and status\\
\midrule
Exact first-promotion law & Original clock rates and exact null elimination; proved here.\\
Positive reduction and tails & Rational identities plus explicit vanishing boundaries; proved here.\\
Rational nonsummability & Pole-orbit argument and polynomial degree comparison; proved here.\\
Primitive rational denominators & Exact prime valuations and classical quadratic Mertens estimates; proved here.\\
Algebraic coefficient heights & Exact prime-ideal valuations, number-field Mertens, and conjugate bounds; proved here.\\
Numerical intervals & Exact finite arithmetic using the analytic tail theorem; reproducible.\\
Global critical transfer & Not constructed.\\
Arithmetic small forms & No suitable signed critical forms or required height/error balance constructed.\\
Transcendence of $\mu$ & Unresolved. The local-family transcendence assertion is also unresolved.\\
\bottomrule
\end{tabular}
\end{center}

\paragraph{Verification.}
The standard-library program \texttt{scripts/check\_first\_promotion.py}
checks independent null-eliminated clock recursions, the finite boundary
identity, tail signs and widths, and good-prime numerator valuations.
Its reference run contains $106{,}865$ exact finite checks.
The JSON report labels floating-point denominator diagnostics as
exploratory. Finite checks do not replace the infinite proofs, establish
asymptotics experimentally, or certify transcendence. The source and exact
reference report are included in the repository, alongside the
original-model and local-model manuscripts.

\paragraph{Acknowledgment.}
AI tools assisted with proof exploration, independent arithmetic audits,
literature retrieval, and preparation of reproducible calculations.

\begin{thebibliography}{9}
\bibitem{KR} P. L. Krapivsky and S. Redner,
\emph{Emergent Network Modularity}, J. Stat. Mech. (2017), 073405.
\url{https://arxiv.org/abs/1706.01514}.
\bibitem{Uzay} C. Uzay, \emph{friend-tree-exponent}, proof revision
\texttt{88f2f9c8c3ebc1b4e0c980db6afd97d29f3df176}.
\url{https://github.com/uzay-eth/friend-tree-exponent}.
\bibitem{Long} C. D. Long, \emph{Effective computation and boundary
robustness of the random friend-tree exponent}, and
\emph{Growth and irrationality in local redirection trees}, research
drafts, November 23, 2026, repository revision
\texttt{6ee9d6e6402699e02860eb02dde319d28a6df385}.
\url{https://github.com/long-mathematics/friend-tree-growth}.
\bibitem{DLMF} NIST Digital Library of Mathematical Functions,
Section 16.2, generalized hypergeometric series.
\url{https://dlmf.nist.gov/16.2}.
\bibitem{Zhao} T. Zhao, \emph{On the mean values of the error terms in
Mertens' theorems}, Research in Number Theory 11 (2025), article 62.
\url{https://doi.org/10.1007/s40993-025-00640-y}.
\bibitem{GarciaLee} S. R. Garcia and E. S. Lee,
\emph{Unconditional explicit Mertens' theorems for number fields and
Dedekind zeta residue bounds}, Ramanujan J. 57 (2022), 1169--1191,
Theorem 1, (A1)--(A2).
\url{https://arxiv.org/abs/2007.10313v4}.
\bibitem{Beukers} F. Beukers, \emph{A refined version of the
Siegel--Shidlovskii theorem}, Ann. Math. 163 (2006), 369--379.
\url{https://annals.math.princeton.edu/wp-content/uploads/annals-v163-n1-p08.pdf}.
\bibitem{Sprindzuk} V. G. Sprind\v zuk,
\emph{The irrationality of the values of certain transcendental
functions}, Math. USSR-Izv. 2 (1968), 89--104, Theorems 1 and 3.
\url{https://doi.org/10.1070/IM1968v002n01ABEH000632}.
\end{thebibliography}
\end{document}
